% !TEX root = ../main.tex
\section{Tổng quan}

\subsection{Giới thiệu chung về Điện toán Đám mây và Serverless Computing}
Gần đây, kiến trúc nguyên khối (Monolithic) đang dần nhường chỗ cho hệ thống phân tán và vi dịch vụ (Microservices). Điện toán đám mây (Cloud Computing) phát triển nhanh chóng, kéo theo nhiều cách thức mới để triển khai ứng dụng. Trong số đó, mô hình Điện toán phi máy chủ (Serverless Computing) đang nhận được nhiều sự chú ý \cite{shahrad2020serverless, hellerstein2018serverless}.

Với Serverless, lập trình viên không còn phải quản lý hay duy trì các máy chủ vật lý hay máy ảo nữa. Họ chỉ cần viết code. Hệ thống sẽ tính tiền dựa trên lượng tài nguyên dùng thực tế (pay-as-you-go). Đặc biệt nhất là tính năng "Scale-to-Zero" – hệ thống có thể tắt toàn bộ các ứng dụng (instances) đang chạy nếu không có request nào, giúp tiết kiệm chi phí đáng kể. Mặc dù rất tiện lợi, nhưng mô hình này lại sinh ra một vấn đề kỹ thuật lớn: lỗi khởi động lạnh.

\subsection{Bài toán Khởi động lạnh (Cold-start Problem)}
Khởi động lạnh (Cold-start) diễn ra khi có một request gửi đến ứng dụng Serverless, nhưng lúc đó ứng dụng lại đang ngủ (nghĩa là số lượng pod bằng 0). Khi đó, hệ thống đám mây sẽ phải làm khá nhiều việc để xử lý request này:
\begin{enumerate}
    \item Xin cấp tài nguyên (CPU, Memory).
    \item Kéo (pull) image từ kho (Registry).
    \item Cấu hình mạng (cấp phát IP).
    \item Chạy ứng dụng bên trong container. Tùy framework mà bước này nhanh hay chậm, như Java thì phải bật JVM, còn Python thì load thư viện.
\end{enumerate}

Tùy vào độ phức tạp của ứng dụng, toàn bộ quá trình này mất khoảng vài giây, có khi lên đến cả chục giây. Hậu quả là request đầu tiên của người dùng sẽ bị nghẽn (block) chờ cho đến khi ứng dụng chạy xong. Đối với những hệ thống đòi hỏi thời gian phản hồi siêu thấp (như giao dịch tài chính hay thời gian thực), độ trễ này ảnh hưởng rất xấu đến trải nghiệm người dùng và dễ gây vi phạm cam kết dịch vụ (SLA).

\subsection{Nền tảng Knative và Hạn chế của Cơ chế mở rộng phản ứng}
Knative là một công cụ mã nguồn mở được chạy trên nền Kubernetes, giúp mang lại tính năng Serverless cho các cụm container \cite{knative2023}. Cơ chế mở rộng số lượng pod mặc định của Knative là Knative Pod Autoscaler (KPA). Nó hoạt động theo kiểu phản ứng (Reactive Autoscaling).

Nói đơn giản là KPA sẽ đợi đến khi thấy lượng truy cập thực tế vượt mức quy định thì mới bắt đầu tăng số pod lên. Mặc dù cách này khá an toàn, nhưng vì tính chất "thấy request mới scale", hệ thống chắc chắn sẽ bị dính khởi động lạnh nếu bỗng nhiên lượng truy cập tăng vọt từ mức 0.

\subsection{Giải pháp Đề xuất: Lập lịch Dự đoán (Predictive Scheduling) với Học máy}
Thay vì thụ động đợi tải tăng, nghiên cứu này xây dựng một giải pháp Lập lịch dự đoán (Predictive Scheduling) chủ động \cite{daw2020predictive}. Mục tiêu là dự đoán trước lượng truy cập để tạo sẵn pod (Pre-warming). Như vậy khi request tới thì pod đã có sẵn rồi, không còn tình trạng khởi động lạnh nữa.

Để làm được việc này, chúng tôi áp dụng mô hình mạng nơ-ron hồi quy dài-ngắn hạn (LSTM) \cite{hochreiter1997long}. LSTM xử lý rất tốt các chuỗi dữ liệu theo thời gian (Time-series data). Dữ liệu về số lượng request (RPS) sẽ được gom lại qua Prometheus, sau đó nạp vào mô hình LSTM để dự báo tải. Dựa trên số liệu dự báo, một chương trình tự viết (Custom Controller) sẽ điều khiển Kubernetes tăng hoặc giảm số pod cho phù hợp.

\subsection{Mục tiêu và Phạm vi nghiên cứu}
\textbf{Mục tiêu chính:}
\begin{itemize}
    \item Phân tích và làm rõ cơ chế hoạt động, cũng như các điểm nghẽn gây ra hiện tượng khởi động lạnh trên nền tảng Knative.
    \item Thiết kế, phát triển và huấn luyện mô hình Học sâu LSTM có khả năng dự đoán biến động lưu lượng mạng với độ chính xác cao.
    \item Phát triển một Kubernetes Custom Controller bằng ngôn ngữ Python có khả năng tương tác với API của Kubernetes để thao tác mở rộng quy mô dựa trên kết quả dự đoán.
    \item Đánh giá hiệu suất của giải pháp đề xuất thông qua các mô phỏng (simulations) về lưu lượng truy cập tăng vọt, so sánh với cơ chế Autoscaler mặc định của Knative.
\end{itemize}

\textbf{Phạm vi nghiên cứu:}
Nghiên cứu này tập trung vào việc tối ưu hóa độ trễ ở cấp độ nền tảng (Platform level) cho các cụm Kubernetes/Knative, không đi sâu vào việc tối ưu hóa mã nguồn (code-level) của từng ngôn ngữ lập trình cụ thể (ví dụ như biên dịch AOT của Java hay tối ưu hóa thư viện Python). Khung đánh giá được giới hạn trong một môi trường mô phỏng cục bộ (Local Kubernetes Cluster) sử dụng các mẫu tải nhân tạo (synthetic workloads) được thiết kế để mô phỏng các đợt tấn công DDoS hoặc sự kiện Flash-sale.

\subsection{Cấu trúc của Luận văn}
Nội dung của luận văn được tổ chức thành 5 chương chính như sau:
\begin{itemize}
    \item \textbf{Chương 1 - Tổng quan:} Giới thiệu bối cảnh, bài toán, giải pháp đề xuất và mục tiêu nghiên cứu.
    \item \textbf{Chương 2 - Cơ sở lý thuyết và Các công nghệ liên quan:} Trình bày chi tiết về kiến trúc Kubernetes, Knative, KPA, và lý thuyết nền tảng về mô hình dự báo chuỗi thời gian LSTM.
    \item \textbf{Chương 3 - Kiến trúc Đề xuất và Hiện thực:} Đi sâu vào thiết kế hệ thống, giải thuật làm mịn (Smoothing Algorithm), cơ chế giả lập tải 3 pha (3-Phase Spike Simulation) và giao diện giám sát (Executive Dashboard).
    \item \textbf{Chương 4 - Đánh giá và Kết quả Thực nghiệm:} Phân tích các số liệu thu thập được từ thực nghiệm, so sánh hiệu năng và đánh giá trực quan qua Dashboard.
    \item \textbf{Chương 5 - Kết luận và Hướng phát triển:} Tóm tắt lại những đóng góp của đề tài và mở ra các định hướng cải tiến trong tương lai.
\end{itemize}